% Generated by scripts/txt2tex.py from chapters/chapter_07.txt. Edit the TXT, not this file.
\chapterquestion{כשהתשובה כבר כתובה על המסך, מי מחליט מה עושים איתה?}
\chapter{האדם במעגל והעתיד}\label{ch:07}
\section{הסמן מפסיק להבהב}
המילה האחרונה הופיעה. הדוגמאות, המשקולות והטוקנים שליווינו עד כאן התכנסו לפסקה מסודרת על בצק. נשאר תפקיד אחד שאף מודל לא ממלא: לקרוא ולהחליט.
הקוראת כבר לא קוראת כמו בתחילת הערב. העצה הראשונה מנוסחת בביטחון רב, והיא נזכרת שביטחון אינו ראיה. היא פותחת מקור נוסף ומשווה, ומנסחת שאלת המשך מדויקת יותר. בלי לשים לב, היא נכנסה לתפקיד שיש לו שם מקצועי: האדם במעגל (human in the loop), החוליה האנושית שמפקחת על תוצר של מערכת אוטומטית, מתקנת אותו ומחליטה אם לפעול לפיו.
\section{גבול שאי אפשר לראות מראש}
למה צריך את החוליה הזאת, אם העוזרים האלה מועילים כל כך? ניסוי שדה עם 758 עובדי ידע, שתוכניתו נרשמה מראש, נערך בשיתוף חברת הייעוץ Boston Consulting Group ופורסם ב-2026, מציע תשובה. במשימות שהיו בתוך גבול היכולת של המערכת, מי שנעזרו בבינה מלאכותית השלימו 12.2\% יותר משימות, סיימו אותן מהר יותר ב-25.1\% והגיעו לאיכות גבוהה יותר. אבל במשימה שהייתה מחוץ לגבול, הסיכוי שלהם להגיע לפתרון נכון היה נמוך ב-19\% לעומת מי שעבדו בלעדיה \cite{S-ch07-04}. החוקרים כינו זאת "הגבול הטכנולוגי המשונן" (jagged technological frontier): אותה מערכת מצטיינת במשימה אחת ומכשילה באחרת, והקו ביניהן אינו נראה מראש \cite{S-ch07-04}.
הספר מציע לחשוב על אפליקציית ניווט. ברוב הנסיעות היא חוסכת זמן ומכירה דרכים שהנהג לא היה מוצא לבד. ואז, יום אחד, היא מכוונת לרחוב שנחסם הבוקר. המסך נשאר רגוע ובטוח בעצמו כתמיד; רק מי שמביט בכביש רואה את המחסום. האדם במעגל הוא הנהג שממשיך להסתכל קדימה.
\begin{labbox}{מה נמדד במקומות עבודה}
שני מחקרים שעברו ביקורת עמיתים, כלומר נבדקו בידי חוקרים אחרים לפני הפרסום, ולכל אחד אוכלוסייה ומשימה משלו.
בניסוי מקוון עם 453 אנשי מקצוע בעלי השכלה אקדמית, שביצעו משימות כתיבה מקצועיות, מי שקיבלו גישה ל-ChatGPT סיימו בממוצע ב-40\% פחות זמן, ואיכות התוצרים עלתה ב-18\%. גם הפערים בין העובדים הצטמצמו \cite{S-ch07-01}.
במוקד תמיכת לקוחות עם 5,172 נציגים, עוזר שיחה מבוסס בינה מלאכותית יוצרת העלה את מספר הפניות שנפתרו בשעה ב-15\% בממוצע. הפחות מנוסים השתפרו גם במהירות וגם באיכות, ואילו אצל המנוסים והמיומנים ביותר נמדדו רווח קטן במהירות וירידה קטנה באיכות \cite{S-ch07-03}.
אף אחד מהמספרים האלה אינו "ההשפעה של AI על העבודה". כל אחד נכון לקבוצה שנבדקה ולמשימה שלה.
\end{labbox}
\section{ללמוד עם הכלי, או במקומו}
ניסוי שדה בתיכון גדול אחד בטורקיה, בסמסטר הסתיו של שנת הלימודים 2023–2024, כלל כמעט 1,000 תלמידי מתמטיקה בכיתות ט' עד י"א, בארבעה מפגשים של 90 דקות \cite{S-ch07-10}. שתי קבוצות עבדו עם GPT-4: אחת בממשק רגיל דמוי ChatGPT, והשנייה בגרסה שהונחתה לתת רמזים בלי למסור את התשובה. בזמן התרגול, כשהעוזר היה זמין, הביצועים עלו ב-48\% וב-127\% בהתאמה לעומת קבוצת הביקורת. אבל במבחן בלי סיוע, הגרסה הפתוחה הורידה את הביצועים ב-17\% ביחס לתלמיד ממוצע בקבוצת הביקורת, ואילו בגרסת הרמזים לא נמצא הבדל מובהק מהביקורת \cite{S-ch07-10}. החוקרים מתארים תלמידים שהשתמשו בגרסה הפתוחה כמי שנשענו עליה כמו על "קב": ביקשו פתרונות והעתיקו אותם. והצ'טבוט עצמו נתן תשובה נכונה רק ב-51\% מהמקרים בממוצע, בתרגילי הניסוי \cite{S-ch07-10}.
מחקר אחר מצייר תמונה הפוכה לכאורה. חוקרים מאוניברסיטת הרווארד ערכו ניסוי מבוקר בקורס מבוא לפיזיקה למדעי החיים, בסתיו 2023, עם 194 סטודנטים: כל קבוצה למדה נושא אחד עם מורה AI בבית ונושא אחר בלמידה פעילה בכיתה. לפי הגרסה שעברה ביקורת עמיתים, עם מורה ה-AI הסטודנטים למדו יותר ובפחות זמן, וההבדל מובהק סטטיסטית, כלומר קשה להסביר אותו כמקרה \cite{S-ch07-13}. אבל המורה הדיגיטלי הזה תוכנן לפי עקרונות פדגוגיים, והמחברים עצמם אינם מניחים שיעלה תמיד על הכיתה, ואינם רואים בו תחליף להוראה פנים אל פנים \cite{S-ch07-13}.
בית ספר אחד וקורס אחד אינם מספיקים כדי לקבוע אם בינה מלאכותית "משפרת" או "פוגעת" בלמידה. יחד הם מראים שעיצוב הכלי ואופן השימוש בו משנים את התוצאה.
\section{מה החוק מבקש מהאדם}
ב-13 ביוני 2024 אומצה באיחוד האירופי תקנה 2024/1689, חוק הבינה המלאכותית (AI Act). ההוראות הכלליות שלה חלות מ-2 בפברואר 2025, והתקנה חלה מ-2 באוגוסט 2026 \cite{S-ch07-09}. סעיף 14, "פיקוח אנושי", דורש שמערכות בסיכון גבוה יתוכננו כך שבני אדם יוכלו לפקח עליהן ביעילות בזמן השימוש, כדי למנוע או לצמצם סיכונים לבריאות, לבטיחות ולזכויות יסוד \cite{S-ch07-08}. המפקחים צריכים להבין את היכולות והמגבלות של המערכת, לפרש נכון את הפלט, להיות מסוגלים להתעלם ממנו, לעקוף אותו או להפוך אותו, ולעצור את המערכת ב"כפתור עצירה" או בהליך דומה \cite{S-ch07-08}. הסעיף חל רק על מה שהתקנה מגדירה כסיכון גבוה, לא על כל שיחה עם צ'טבוט, אבל רשימת היכולות שלו מתאימה גם לערב במטבח.
\begin{conceptbox}{הטיית אוטומציה (automation bias)}
בסעיף 14 מופיע מונח שנותן שם לפיתוי לוותר על הבדיקה: הטיית אוטומציה, הנטייה להסתמך באופן אוטומטי, או יותר מדי, על הפלט של מערכת \cite{S-ch07-08}. החוק אינו מסתפק בכך שהמפקח יכול לעקוף את המערכת; הוא דורש שיהיה מודע לנטייה הזאת. גם חומר הלימוד של גוגל על למידת מכונה מונה אותה בין סוגי ההטיה \cite{S-ch02-14}. כאן היא מוצגת כמושג מוגדר, לא כממצא על מידת השכיחות שלה.
\end{conceptbox}
\section{ארגז הכלים של הבודקת}
מערכות עם רכיבי בינה מלאכותית פועלות בתנאי אי-ודאות. לפעמים הן מתריעות על משהו שאינו קיים, ולפעמים מחמיצות משהו שקיים (false positives ו-false negatives), ולכן התנהגותן עלולה להיות בלתי צפויה, מבלבלת ואף מסוכנת \cite{S-ch07-06}. סלימה אמרשי (Saleema Amershi) ועמיתיה ריכזו יותר מ-150 המלצות עיצוב ל-18 הנחיות לאינטראקציה בין אדם למערכת AI, ובחנו אותן עם 49 מעצבים מול 20 מוצרים נפוצים \cite{S-ch07-05,S-ch07-06}. ההנחיות נכתבו למעצבים, אבל כאן נהפוך אותן לשאלות של משתמשת: מה המערכת יודעת לעשות (G1)? באיזו תדירות היא טועה (G2)? למה ענתה כך (G11)? וכמה קל לתקן אותה כשהיא שוגה (G9) \cite{S-ch07-06}?
לפי אונסק"ו, מוסדות החינוך עדיין אינם ערוכים במידה רבה לאמת את הכלים, והארגון ממליץ שבתהליכי התיקוף האתי והעיצוב הפדגוגי האדם יישאר הגורם הפועל \cite{S-ch07-07}.
\begin{figure}[htbp]
\centering
% Figure 7.1: the reader's decision cycle. The four questions adapt guidelines G1, G2, G11, G9 (S-ch07-06).
\begin{bookdiagram}[x=1cm,y=1cm]
  \tikzset{nd/.style={blk,draw=turquoise,line width=1.1pt,text width=2.5cm,minimum height=10mm}}
  \def\R{2.9}
  \node[nd] (n1) at (90:\R)   {שואלת};
  \node[nd] (n2) at (18:\R)   {מקבלת תשובה};
  \node[nd,draw=violet,line width=1.6pt,fill=violet!8] (n3) at (-54:\R) {בודקת מול מקורות};
  \node[nd] (n4) at (-126:\R) {מחליטה};
  \node[nd] (n5) at (162:\R)  {מחדדת את השאלה};
  \foreach \a/\b in {n1/n2,n2/n3,n3/n4,n4/n5,n5/n1}{\draw[-{Stealth[length=2.6mm]},line width=1pt,midnight] (\a) to[bend left=18] (\b);}
  \node[draw=gold,line width=1.6pt,rounded corners=2mm,fill=white,font=\headingfont,inner sep=2.5mm] at (0,0) {האדם במעגל};
  \node[draw=violet,line width=0.8pt,rounded corners=2mm,fill=white,text width=3.3cm,align=right,font=\footnotesize,anchor=west]
     (q) at (4.0,-2.6) {{\headingfont ארבע שאלות לבדיקה}\par מה המערכת יודעת?\par כמה היא טועה?\par למה ענתה כך?\par איך מתקנים?};
  \draw[violet,line width=0.8pt,dashed] (n3.east) -- (q.west);
\end{bookdiagram}

\caption{מעגל ההחלטה של הקוראת}\label{fig:7.1}
\figurenote{מקור: מסגרת הספר; ארבע השאלות הן עיבוד של הספר להנחיות G1, G2, G11 ו-G9 ב-\cite{S-ch07-06}}
\end{figure}
\begin{thinkbox}
התלמידים שעבדו עם הגרסה הפתוחה בניסוי המתמטיקה לא הרגישו שביצעו פחות טוב או שלמדו פחות, אף שבמבחן בלי סיוע הם השיגו פחות \cite{S-ch07-10}. אם התחושה הפנימית אינה מדד אמין, איך תדעו שלמדתם משהו בעזרת צ'טבוט, ולא רק העתקתם? באילו משימות שלכם הוא נמצא, לדעתכם, מחוץ לגבול? ומי אחראי כשהעוזר שוגה, וההחלטה נשארה בידיכם?
\end{thinkbox}
\section{מה תפח בסוף}
הספר לא יגלה אם הבצק של הקוראת תפח. זו שאלה של מטבח, לא של בינה מלאכותית. אבל משהו אחר גדל לאורך הערב. היא התחילה עם שאלה ועם אמון; היא מסיימת עם שאלה ועם שיטה: לשאול, להשוות, להחליט, ולשאול שוב, טוב יותר. מה שתפח בסוף הוא שיקול הדעת שלה.
ומה לגבי העתיד? תחזיות על משרות שייעלמו או ייווצרו קל לנסח וקשה לבסס, ולכן אין כאן כאלה. המחקרים בפרק הזה מלמדים משהו צנוע יותר: התוצאה תלויה במשימה, באדם ובאופן שבו עוצבה, והגבול בין הצלחה לכישלון אינו מסומן. מי שיודעים לבדוק יוכלו לנוע לאורכו.
ונשאר מעגל אחד לסגור. הספר שבידיכם נכתב בעזרת סוכני בינה מלאכותית: אחד אימת מקורות, אחד כתב, אחד ביקר והחזיר לתיקון, וסוכן-על תיאם ביניהם. מעל כולם עמד אדם שאישר או דחה. בנספח נפתח גם את הקלעים האלה.
\begin{insightsbox}
\begin{itemize}
\item הצלחה במשימה אחת אינה ערובה להצלחה בבאה, ולכן בודקים כל תשובה מחדש \cite{S-ch07-04}.
\item כל נתון על השפעת AI תקף רק לאוכלוסייה ולמשימה שנבדקו; אנשי מקצוע במשימות כתיבה ונציגי תמיכה אינם "כל העובדים" \cite{S-ch07-01,S-ch07-03}.
\item בניסוי בתיכון אחד, שיפור בזמן התרגול לא הבטיח למידה, ובגרסה שנתנה רמזים במקום תשובות לא נמצא נזק מובהק \cite{S-ch07-10}.
\item פיקוח אנושי מעוגן בחוק האירופי למערכות בסיכון גבוה: להבין, לפרש, לעקוף ולעצור, ולהכיר את הפיתוי לסמוך יותר מדי \cite{S-ch07-08}.
\item ארבע שאלות פשוטות, מה היא יודעת, כמה היא טועה, למה ענתה כך ואיך מתקנים, הופכות קוראת לבודקת \cite{S-ch07-06}.
\end{itemize}
\end{insightsbox}
